\begin{figure}[H]
        \centering

        \begin{subfigure}[t]{0.3\textwidth}
            \centering
            \includesvg[width=\textwidth]{img/graphs/6/F3/1}
            \caption{
                Plocha S = (201,09 ± 201,14) $\mathrm{mN \cdot s}$
            }
            \label{fig:integrales-3a}
        \end{subfigure}%
    \hfill
\begin{subfigure}[t]{0.3\textwidth}
            \centering
            \includesvg[width=\textwidth]{img/graphs/6/F3/2}
            \caption{
                Plocha S = (197,9 ± 196,1) $\mathrm{mN \cdot s}$
            }
            \label{fig:integrales-3b}
        \end{subfigure}%
    \hfill
\begin{subfigure}[t]{0.3\textwidth}
            \centering
            \includesvg[width=\textwidth]{img/graphs/6/F3/3}
            \caption{
                Plocha S = (195,8 ± 203,4) $\mathrm{mN \cdot s}$
            }
            \label{fig:integrales-3c}
        \end{subfigure}

    \vspace{0.5em}
\begin{subfigure}[t]{0.3\textwidth}
            \centering
            \includesvg[width=\textwidth]{img/graphs/6/F3/4}
            \caption{
                Plocha S = (162,6 ± 168,3) $\mathrm{mN \cdot s}$
            }
            \label{fig:integrales-3d}
        \end{subfigure}%
    \hfill
\begin{subfigure}[t]{0.3\textwidth}
            \centering
            \includesvg[width=\textwidth]{img/graphs/6/F3/5}
            \caption{
                Plocha S = (189,4 ± 195,5) $\mathrm{mN \cdot s}$
            }
            \label{fig:integrales-3e}
        \end{subfigure}%
    \hfill
\begin{subfigure}[t]{0.3\textwidth}
            \centering
            \includesvg[width=\textwidth]{img/graphs/6/F3/6}
            \caption{
                Plocha S = (178,7 ± 179,9) $\mathrm{mN \cdot s}$
            }
            \label{fig:integrales-3f}
        \end{subfigure}

    \vspace{0.5em}
\begin{subfigure}[t]{0.3\textwidth}
            \centering
            \includesvg[width=\textwidth]{img/graphs/6/F3/7}
            \caption{
                Plocha S = (185,8 ± 191,8) $\mathrm{mN \cdot s}$
            }
            \label{fig:integrales-3g}
        \end{subfigure}%
    \hfill
\begin{subfigure}[t]{0.3\textwidth}
            \centering
            \includesvg[width=\textwidth]{img/graphs/6/F3/8}
            \caption{
                Plocha S = (184,7 ± 187,8) $\mathrm{mN \cdot s}$
            }
            \label{fig:integrales-3h}
        \end{subfigure}%
    \hfill
\begin{subfigure}[t]{0.3\textwidth}
            \centering
            \includesvg[width=\textwidth]{img/graphs/6/F3/9}
            \caption{
                Plocha S = (188,0 ± 185,6) $\mathrm{mN \cdot s}$
            }
            \label{fig:integrales-3i}
        \end{subfigure}

    \vspace{0.5em}
\begin{subfigure}[t]{0.3\textwidth}
            \centering
            \includesvg[width=\textwidth]{img/graphs/6/F3/10}
            \caption{
                Plocha S = (199,2 ± 206,8) $\mathrm{mN \cdot s}$
            }
            \label{fig:integrales-3j}
        \end{subfigure}

        \caption{
            Naměřené hodnoty síly $F$ v závislosti na čase $t$. Výpočet plochy pod integrálem. Vypočteno skrze program Python \cite{GitHubMAT}.
        }
        \label{fig:integrales-3}
    \end{figure}
    